% \makeatletter
% \renewcommand{\tikz@align@newline@}{%  The default definition had 0pt here VVV
%   \unskip \pgfutil@ifnextchar [\tikz@@align@newline {\tikz@@align@newline [5pt]}%
% }
% % For nodes with text width
% \renewcommand{\@xcentercr}[1][5pt]{%
%   \addvspace{-\parskip}\vskip#1\ignorespaces
% }
% \makeatother

\tikzset{
  blocknode/.style={thick,  
    text width=0.95\linewidth},
  defnode/.style={blocknode,draw=color3 ,fill=color3!5!white,
  },
  thmnode/.style={blocknode,draw=color2, fill=color2!5!white}
}

\renewenvironment{theorem}[1][]{%
  \addtobeamertemplate{itemize/enumerate body begin}{}{\vspace{0.5em}}
  \begin{center}
  \begin{tikzpicture}
    \node[thmnode] (thm) \bgroup
      \textcolor{color2}{\textbf{Theorem} #1}\\
    }{%
      \egroup;
      %\draw[thick,color2] (thm.north west) -- (thm.south west);
    \end{tikzpicture}
  \end{center}
}

\renewenvironment{definition}[1][Definition]{%
  \addtobeamertemplate{itemize/enumerate body begin}{}{\vspace{0.5em}}
  %\setbeamercolor{alerted text}{fg=color3}
  \begin{center}
  \begin{tikzpicture}
    \node[defnode] (def) \bgroup
      \textcolor{color3}{\textbf{#1}}\\
    }{
      \egroup;
      % \draw[thick,color2] (def.south west) -- (def.north west) % --++ (2,0)
      % ;
    \end{tikzpicture}
  \end{center}
}


%%% Local Variables:
%%% mode: latex
%%% TeX-master: t
%%% End:
